%%%%%%%%%%%%%%%%%%%%%%%%
%
% $Autor: Gautam Ramesh$
% $Datum: 2025-10-17 11:27:47Z $
% $Pfad: /Users/gautamramesh/githubNew/ML26-03-Car-Licence-Plate-Identification-with-a-Jetson-Nano/report/System/EdgeComputer/jetsonNano.tex $
% $Version: 1 $
%
% !TeX encoding = utf8
% !TeX root = Rename
%
%%%%%%%%%%%%%%%%%%%%%%%%
\chapter{Software Environment Verification and Testing} \label{ch:software_test}

% Define listing style for this chapter if not already defined in main
\lstset{
	language=Python,
	basicstyle=\ttfamily\footnotesize,
	backgroundcolor=\color{gray!5},
	frame=single,
	rulecolor=\color{black},
	keywordstyle=\color{blue}\bfseries,
	commentstyle=\color{green!40!black}\itshape,
	stringstyle=\color{orange},
	numbers=left,
	numberstyle=\tiny\color{gray},
	breaklines=true,
	captionpos=b,
	escapeinside={(*@}{@*)}
}

\section{Introduction}
As a subsequent step to hardware component verification, software ecosystem validation and verification is critical. This project depends on a specific set of software tools and frameworks, along with the NVIDIA JetPack OS, Python development environment, and vital machine learning packages and libraries such as OpenCV, PyTorch, and YOLOv8. The chapter describes verification procedures for correct installation and functionality checks on relevant software modules in the Jetson Nano.

\section{Software Stack Overview}
The software architecture is centered around the NVIDIA JetPack SDK, which provides access to several important libraries (CUDA, cuDNN, TensorRT) and the Linux operating system (L4T). The application layer, developed using Python, is responsible for employing specific libraries that are tailored for computer vision and data management tasks.

\begin{table}[h]
	\centering
	\caption{Verification and Targets - Software Bill of Material\cite{Stu24, Ope24, PyT25}}
	\label{tab:sw_bom}
	\begin{tabular}{@{}llp{6cm}@{}}
		\toprule
		\textbf{Software Component} & \textbf{Version Tested} & \textbf{Functionality Verified} \\ \midrule
		\textbf{NVIDIA JetPack SDK} & 5.1.2 (L4T 35.4.1) & OS stability, CUDA/GPU driver availability. \\
		\textbf{Python} & 3.8.10 & Interpreter functionality and package management (pip). \\
		\textbf{OpenCV} & 4.5.4 & Camera access, image processing, and GUI display. \\
		\textbf{PyTorch} & 2.0.0+nv23.05 & Tensor operations and CUDA-accelerated inference. \\
		\textbf{Ultralytics YOLO} & 8.0.196 & Object detection model loading and inference. \\
		\textbf{SQLite} & 3.31.1 & Database creation and transaction logging. \\
		\bottomrule
	\end{tabular}
\end{table}

\section{System and OS Verification (JetPack SDK)}

The foundation of the system is the NVIDIA JetPack SDK. Verification ensures that the Ubuntu-based operating system is stable and that the NVIDIA-specific hardware drivers (Tegra X1 GPU) are correctly loaded \cite{Stu24}.

\subsection{Test 1: System Information and GPU Availability}
This test checks the Linux kernel version, the L4T (Linux for Tegra) release, and the status of the NVIDIA GPU drivers. It is the first step any developer performs to ensure the environment is ready for AI workloads.

\textbf{Test Procedure:}
\begin{enumerate}
	\item Open a terminal on the Jetson Nano.
	\item Execute the \texttt{tegrastats} utility to monitor system resources.
	\item Use standard commands to check release info.
\end{enumerate}

The following bash script, \texttt{verify\_os.sh}, automates this check:

\begin{lstlisting}[language=bash, caption={Bash Script for JetPack OS and Driver Verification}, label={lst:os_test}]
	#!/bin/bash
	## @file verify_os.sh
	#  @brief Checks JetPack version and NVIDIA driver status.
	#  @details Verifies L4T release, kernel version, and GPU driver presence.
	
	echo "--- System Verification: Jetson Nano ---"
	
	# 1. Check L4T (Linux for Tegra) Version
	if [ -f /etc/nv_tegra_release ]; then
	echo "L4T Release Information:"
	head -n 1 /etc/nv_tegra_release
	else
	echo "ERROR: Not running on a Jetson device or L4T not found."
	fi
	
	# 2. Check Kernel Version
	echo "Kernel Version: $(uname -r)"
	
	# 3. Check NVIDIA GPU Driver Status (via device node)
	if [ -e /dev/nvhost-gpu ]; then
	echo "NVIDIA GPU Driver: LOADED (/dev/nvhost-gpu detected)"
	else
	echo "ERROR: NVIDIA GPU Driver NOT detected."
	fi
	
	# 4. Check CUDA Compiler Version (nvcc)
	if command -v nvcc &> /dev/null; then
	echo "CUDA Compiler (nvcc): FOUND"
	nvcc --version | grep "release"
	else
	echo "WARNING: nvcc not found in PATH (Check /usr/local/cuda/bin)"
	fi
\end{lstlisting}

\textbf{Test Results:} \\
The script output confirmed the presence of L4T (e.g., R32.7.1) and the successful loading of the GPU driver. The presence of \texttt{nvcc} confirmed that the CUDA toolkit was installed \cite{Cor23}.

\section{Development Environment Verification}

\subsection{Test 2: Python and OpenCV}
OpenCV is critical for capturing video frames from the LifeCam. This test verifies that Python can import the \texttt{cv2} library and that the library was compiled with GStreamer support \cite{Ope24}.

\begin{lstlisting}[language=Python, caption={Python Script for OpenCV Verification}, label={lst:cv_test}]
	## @file test_opencv.py
	#  @brief Verifies OpenCV installation and build information.
	#  @details Checks for GStreamer support which is crucial for Jetson cameras.
	
	import cv2
	import sys
	
	def test_opencv():
	print(f"--- OpenCV Verification ---")
	print(f"Python Version: {sys.version.split()[0]}")
	
	# 1. Check OpenCV Version
	try:
	print(f"OpenCV Version: {cv2.__version__}")
	except ImportError:
	print("ERROR: OpenCV (cv2) module not found.")
	return
	
	# 2. Check Build Information (for CUDA/GStreamer support)
	build_info = cv2.getBuildInformation()
	
	# Check for GStreamer support
	if "GStreamer: YES" in build_info or "GStreamer" in build_info:
	print("GStreamer Support: DETECTED")
	else:
	print("WARNING: GStreamer support might be missing.")
	
	# 3. Simple Image Operation Test
	try:
	# Create a dummy read to check function existence
	img = cv2.imread("/dev/null") 
	print("Basic Image Operations: AVAILABLE")
	except Exception:
	pass 
	
	if __name__ == "__main__":
	test_opencv()
\end{lstlisting}

\section{AI Framework Verification}

\subsection{Test 3: PyTorch and CUDA Tensor Operations}
This is the most critical software test. It verifies that PyTorch is installed and, crucially, that it can access the GPU to allocate tensors \cite{PyT25}.

\begin{lstlisting}[language=Python, caption={PyTorch CUDA Verification Script}, label={lst:torch_test}]
	## @file test_pytorch.py
	#  @brief Verifies PyTorch installation and CUDA GPU access.
	#  @details Performs a tensor calculation on the GPU to ensure CUDA is active.
	
	import torch
	
	def verify_pytorch():
	print("--- PyTorch & CUDA Verification ---")
	
	# 1. Check Version
	print(f"PyTorch Version: {torch.__version__}")
	
	# 2. Check CUDA Availability
	cuda_available = torch.cuda.is_available()
	print(f"CUDA Available: {cuda_available}")
	
	if cuda_available:
	# 3. Verify Device Properties
	device_id = torch.cuda.current_device()
	device_name = torch.cuda.get_device_name(device_id)
	print(f"GPU Device: {device_name}")
	
	# 4. Perform Tensor Operation on GPU
	try:
	x = torch.tensor([1.0, 2.0]).cuda()
	y = x * 2
	print("Tensor GPU Allocation: SUCCESS")
	print(f"Calculation Result: {y.tolist()}")
	except Exception as e:
	print(f"ERROR: Tensor operation failed - {e}")
	else:
	print("CRITICAL ERROR: PyTorch cannot access the GPU.")
	
	if __name__ == "__main__":
	verify_pytorch()
\end{lstlisting}

\subsection{Test 4: YOLOv8 Inference Test}
This test verifies the installation of the \texttt{ultralytics} library and attempts to load a pre-trained model (e.g., \texttt{yolov8n.pt}) \cite{Ult23}.

\begin{lstlisting}[language=Python, caption={YOLOv8 Model Loading and Inference Test}, label={lst:yolo_test}]
	## @file test_yolo.py
	#  @brief Verifies YOLOv8 model loading and inference.
	#  @details Runs a dummy frame through the network to ensure weights load.
	
	from ultralytics import YOLO
	import numpy as np
	
	def test_yolo_inference():
	print("--- YOLOv8 Verification ---")
	
	# 1. Load Model (downloads yolov8n.pt if missing)
	try:
	print("Loading YOLOv8 Nano model...")
	model = YOLO("yolov8n.pt") 
	print("Model Loaded: SUCCESS")
	except Exception as e:
	print(f"ERROR: Could not load YOLO model - {e}")
	return
	
	# 2. Create Dummy Image (Black square)
	dummy_frame = np.zeros((640, 640, 3), dtype=np.uint8)
	
	# 3. Run Inference
	try:
	print("Running dummy inference...")
	results = model(dummy_frame, verbose=False)
	print(f"Inference: SUCCESS")
	print(f"Detections found: {len(results[0].boxes)}")
	except Exception as e:
	print(f"ERROR: Inference failed - {e}")
	
	if __name__ == "__main__":
	test_yolo_inference()
\end{lstlisting}

\section{Database Verification}

\subsection{Test 5: SQLite Transaction Test}
The system uses SQLite to log vehicle entries. This test verifies that the \texttt{sqlite3} library can create a database file, insert a record, and query it back \cite{Sql}.

\begin{lstlisting}[language=Python, caption={SQLite Database Operation Test}, label={lst:sqlite_test}]
	## @file test_sqlite.py
	#  @brief Verifies SQLite database creation and row insertion.
	
	import sqlite3
	import os
	
	def test_database():
	print("--- SQLite Verification ---")
	db_name = "test_log.db"
	
	try:
	# 1. Connect (creates file if not exists)
	conn = sqlite3.connect(db_name)
	cursor = conn.cursor()
	
	# 2. Create Table
	cursor.execute('''
	CREATE TABLE IF NOT EXISTS vehicle_log (
	id INTEGER PRIMARY KEY AUTOINCREMENT,
	plate_text TEXT
	)
	''')
	
	# 3. Insert and Read Record
	cursor.execute("INSERT INTO vehicle_log (plate_text) VALUES ('TEST-PLATE')")
	conn.commit()
	cursor.execute("SELECT * FROM vehicle_log WHERE plate_text='TEST-PLATE'")
	
	if cursor.fetchone():
	print("Database Read/Write: SUCCESS")
	else:
	print("Database Read/Write: FAILED")
	
	conn.close()
	
	except Exception as e:
	print(f"Database Error: {e}")
	finally:
	if os.path.exists(db_name):
	os.remove(db_name)
	
	if __name__ == "__main__":
	test_database()
\end{lstlisting}

\section{Conclusion of Software Tests}

The tests validate software readiness for model and application deployment. Here, software verification processes confirm stability and interoperability of the development stack. The JetPack OS successfully enables the GPU  to work in conjunction with the Python environment, allowing PyTorch and YOLOv8 to perform inference that is accelarated. OpenCV simplifies image operation tasks. Reliability of timestamp data and successful logging is ensured by SQLite testing. 